\section{Introduction}
\label{sec:introduction}

Autonomous driving (AD) systems have advanced remarkably over the past decade, progressing
from rule-based lane-keeping assistants to sophisticated neural architectures capable of
navigating complex urban environments~\cite{caesar2020nuscenes,sun2020scalability,
geiger2012we}. Yet a fundamental limitation persists: contemporary end-to-end systems
remain opaque. They cannot articulate \emph{why} they braked at a particular intersection,
\emph{how} they anticipated a pedestrian's trajectory, or \emph{what} would have occurred had
a different manoeuvre been chosen. This lack of interpretability erodes public trust,
complicates regulatory approval, and makes post-incident analysis arduous.

Large language models (LLMs) offer a compelling avenue toward explainability. Equipped with
broad world knowledge, rich commonsense priors, and the ability to generate fluent natural
language, models such as GPT-4~\cite{openai2023gpt4}, LLaMA~\cite{touvron2023llama2}, and
Gemini~\cite{google2023gemini} have demonstrated impressive performance on scene understanding,
question answering, and instruction-following tasks. When integrated with perception backbones
via vision-language interfaces~\cite{radford2021learning,li2023blip2,liu2023visual},
these models can reason over visual scenes and generate human-readable explanations of
driving decisions.

\subsection*{The Causal Reasoning Gap}

Despite rapid progress, the \emph{data infrastructure} required to train and rigorously
benchmark LLM-based driving agents for causal and counterfactual reasoning remains immature.
Existing language-annotated driving datasets can be broadly divided into two categories.
\emph{Perception-centric} datasets (\eg{} NuScenes-QA~\cite{qian2024nuscenes_qa},
STRIDE-QA~\cite{wu2024stride}, Box-QAymo~\cite{boxqaymo2024}) annotate object existence,
spatial relations, and attribute recognition but do not model the \emph{causal structure}
of driving events.
\emph{Explanation-centric} datasets (\eg{} BDD-X~\cite{kim2018textual},
HAD~\cite{chen2019had}, DRAMA~\cite{malla2023drama}) provide post-hoc rationale for
driving decisions but are limited in scale, geographic diversity, and the depth of
counterfactual exploration.

No existing benchmark answers questions of the form:
\begin{enumerate}[leftmargin=2em, noitemsep]
  \item \emph{``Why did the vehicle decelerate sharply at time $t$?''} (causal QA)
  \item \emph{``What would have happened if the cyclist had not swerved left?''} (counterfactual QA)
  \item \emph{``Given this icy road and oncoming truck, what is the safest sequence of manoeuvres?''} (risk-aware planning)
\end{enumerate}
Answering such questions requires datasets annotated with event-level causal chains,
counterfactual premises and outcomes, and multi-agent intention models --- none of which
are available at scale in the existing literature.

\subsection*{The \lucid{} Dataset}

We introduce \lucid{} (\textbf{L}anguage-grounded \textbf{U}nderstanding for
\textbf{C}ausal \textbf{I}nference and \textbf{D}ecision-making in autonomous
\textbf{D}riving), the first dataset specifically designed to support causal and
counterfactual reasoning for LLM-based autonomous driving. \lucid{} was constructed
through an 18-month multi-city collection campaign and a rigorous three-stage human
annotation protocol.

The key properties of \lucid{} are as follows:
\begin{itemize}[leftmargin=2em, noitemsep]
  \item \textbf{Scale}: 12,847 driving scenarios totalling $\sim$107 hours of driving,
        yielding 2,312,460 camera frames at 10 FPS across eight synchronized surround-view cameras.
  \item \textbf{Language}: 847,293 language annotations spanning five types --- causal chain QA,
        counterfactual QA, planning instructions, scene descriptions, and risk assessments.
  \item \textbf{Geographic diversity}: Six cities across three continents (Shanghai, Beijing,
        Singapore, London, San Francisco, Munich), capturing heterogeneous traffic
        cultures and infrastructure.
  \item \textbf{Condition diversity}: Seven weather states (clear, cloudy, rainy, snowy,
        foggy, hazy, mixed) and four lighting states (daytime, dusk/dawn, night, low-light adverse).
  \item \textbf{Sensor richness}: Eight cameras, five LiDAR units, five mmWave radars,
        GPS/IMU, and lane-level HD maps with proprietary semantic annotations.
  \item \textbf{Novelty}: The first dataset with \emph{structured counterfactual QA pairs}
        and \emph{multi-step causal chain annotations} at scale.
\end{itemize}

\subsection*{Contributions}

This paper makes the following contributions:

\begin{enumerate}[leftmargin=2em, noitemsep, label=\textbf{C\arabic*.}]
  \item We design and release \lucid{}, the first large-scale driving dataset
        with structured causal and counterfactual reasoning annotations
        (312,456 causal QA pairs and 198,327 counterfactual QA pairs).
  \item We propose a hierarchical annotation schema that captures driving events at
        the scene, agent, event, and decision levels, enabling fine-grained evaluation
        of reasoning depth.
  \item We introduce four benchmarks tasks --- CSU, CFR, IFN, and RAP --- each with
        well-defined evaluation protocols and automated metrics.
  \item We conduct a comprehensive evaluation of six state-of-the-art vision-language
        models, establishing competitive baselines and identifying key failure modes.
  \item We publicly release all data, annotation tools, evaluation code, and model
        checkpoints to facilitate reproducible research.
\end{enumerate}

\subsection*{Paper Organisation}

The remainder of this paper is organised as follows. Section~\ref{sec:related_work}
reviews related datasets and methods. Section~\ref{sec:design} describes the design
objectives and annotation schema of \lucid{}. Section~\ref{sec:collection} details
the data collection methodology. Section~\ref{sec:annotation} describes the
annotation pipeline and quality control. Section~\ref{sec:statistics} presents
dataset statistics. Section~\ref{sec:benchmarks} defines the four benchmark tasks.
Section~\ref{sec:experiments} reports baseline experiments. Section~\ref{sec:comparison}
compares \lucid{} with existing datasets. Section~\ref{sec:limitations} discusses
limitations and future directions, and Section~\ref{sec:conclusion} concludes.
